% Assume-se que \pretextual já foi feito

\imprimircapa%
\imprimirfolhaderosto*
% \protect\incluirfichacatalografica{ficha.pdf}  % descomentar quando tiver a ficha
{}% placeholder: evita que \imprimirfolhadecertificacao seja engolido pelo \@ifstar
\imprimirfolhadecertificacao


% \begin{dedicatoria}
% \end{dedicatoria}


\begin{agradecimentos}
Ao Prof.\ Márcio Bastos Castro, pela orientação cuidadosa, pela liberdade de
exploração e pelo ambiente de pesquisa que tornou este trabalho possível.

Ao LaPeSD, em especial a Vanderlei Munhoz, cujo desenvolvimento do orquestrador
HPC@Cloud serviu de base para a integração proposta neste trabalho, e cujas
contribuições técnicas ao longo do projeto foram determinantes.

Este trabalho foi parcialmente financiado pelo Conselho Nacional de
Desenvolvimento Científico e Tecnológico (CNPq) e pela Amazon Web Services
(AWS) por meio da Chamada CNPq/AWS Nº~64/2022 (Cloud Credits for Research).

À minha família, pelo apoio incondicional ao longo de toda a graduação.
\end{agradecimentos}


% \begin{epigrafe}
% \end{epigrafe}


\begin{resumo}[Resumo]
Instâncias \textit{spot} da Amazon Web Services oferecem descontos de até 90\%
em relação ao preço convencional, mas podem ser revogadas pelo provedor a
qualquer momento, com aviso de apenas dois minutos, descartando todo o progresso
computacional acumulado em aplicações MPI de longa duração. Este trabalho propõe
e avalia uma arquitetura de execução resiliente para aplicações MPI legadas em
instâncias \textit{spot}, sem exigir modificação de código, por meio da
integração do MANA, sistema de \textit{checkpoint/restart} transparente para MPI,
ao orquestrador HPC@Cloud. A topologia adotada é híbrida: nó coordenador em
instância \textit{burstable on-demand} e nós de computação em instâncias
\textit{spot}.
Duas estratégias de recuperação são implementadas: a
\textsc{Replace}, que reprovísiona o nó interrompido e retoma a execução com a
capacidade original do \textit{cluster}, e a \textsc{Degraded}, que marca o nó
como inativo e reinicia imediatamente com $N{-}1$ \textit{workers}. Uma terceira
opção, a \textsc{Adaptive}, escolhe automaticamente entre as duas no momento da
falha a partir de uma heurística calibrada empiricamente, validada tanto nos
dados de treino quanto em uma carga de trabalho sintética fora da amostra,
mantendo a mesma decisão independentemente do tamanho do \textit{cluster}. A
avaliação
experimental, conduzida com 282 execuções sobre os \textit{benchmarks} NPB
CG-C, EP-D e LU-C em \textit{clusters} de 2, 4 e 8 \textit{workers}, revelou
dois regimes de viabilidade: para \textit{jobs} curtos como o CG-C (12 a 35~s),
o \textit{overhead} de recuperação domina e a abordagem não se justifica; para
\textit{jobs} de duração moderada a longa como o EP-D e o LU-C, a
\textsc{Degraded} alcançou menor tempo de execução em 17 das 18 configurações
avaliadas, com economia de custo de 2 a 38\% em relação à execução equivalente
em instâncias \textit{on-demand} sem tolerância a falhas.

\vspace{\baselineskip}
\textbf{Palavras-chave:} Computação em Nuvem. Computação de Alto Desempenho.
Instâncias \textit{Spot}. Tolerância a Falhas. \textit{Checkpoint/Restart}.
MANA. HPC@Cloud.
\end{resumo}


\begin{abstract}
AWS spot instances offer discounts of up to 90\% compared to on-demand pricing,
but may be revoked by the provider at any time with only two minutes of notice,
discarding all accumulated computational progress in long-running MPI
applications. This work proposes and evaluates a resilient execution
architecture for legacy MPI applications on AWS spot instances, without
requiring code modification, by integrating MANA, a transparent
\textit{checkpoint/restart} system for MPI, with the HPC@Cloud orchestrator.
The adopted topology is hybrid: a coordinator node on a \textit{burstable
on-demand} instance and compute nodes on \textit{spot} instances.
Two recovery
strategies are implemented: \textsc{Replace}, which reprovisions the interrupted
node and resumes execution with the original cluster capacity, and
\textsc{Degraded}, which marks the node as inactive and immediately restarts
with $N{-}1$ workers. A third option, \textsc{Adaptive}, automatically chooses
between the two at the moment of failure based on an empirically calibrated
heuristic, validated both on the training data and on a synthetic workload
outside that sample, keeping the same decision regardless of cluster size. The
experimental evaluation, conducted with 282 executions
on the NPB benchmarks CG-C, EP-D, and LU-C on clusters of 2, 4, and 8 workers,
revealed two viability regimes: for short jobs such as CG-C (12 to 35~s),
recovery overhead dominates and the approach is not worthwhile; for medium to
long-running jobs such as EP-D and LU-C, \textsc{Degraded} achieved shorter
execution times in 17 out of 18 evaluated configurations, with cost savings of
2 to 38\% compared to the equivalent on-demand execution without fault
tolerance.

\vspace{\baselineskip}
\textbf{Keywords:} Cloud Computing. High Performance Computing. Spot Instances.
Fault Tolerance. Checkpoint/Restart. MANA. HPC@Cloud.
\end{abstract}


\listoffigures*
\listoftables*
\listofalgorithms*

%%% Força todas as siglas definidas em acronyms.tex a aparecer na lista,
%%% mesmo sem uso de \gls{} no texto.
\glsaddall
\listadesiglas*[7em]

\tableofcontents*%

%%% Local Variables:
%%% mode: latex
%%% TeX-master: "main"
%%% End:
